% Created 2026-06-16 mar 23:57
% Intended LaTeX compiler: pdflatex
\documentclass[11pt]{article}
\usepackage[utf8]{inputenc}
\usepackage[T1]{fontenc}
\usepackage{graphicx}
\usepackage{longtable}
\usepackage{wrapfig}
\usepackage{rotating}
\usepackage[normalem]{ulem}
\usepackage{amsmath}
\usepackage{amssymb}
\usepackage{capt-of}
\usepackage{hyperref}
\author{Abraham Nuñez}
\date{\today}
\title{Seguridad en Redes de Computo - Notas Semestre VI}
\hypersetup{
 pdfauthor={Abraham Nuñez},
 pdftitle={Seguridad en Redes de Computo - Notas Semestre VI},
 pdfkeywords={},
 pdfsubject={},
 pdfcreator={Emacs 30.1 (Org mode 9.7.39)}, 
 pdflang={English}}
\begin{document}

\maketitle
\tableofcontents

\section{Introducción}
\label{sec:org86b12c2}
\begin{itemize}
\item Profesor: Martha Elizabet Dominguez Barcenas
\item Marteria: Seguridad en Redes de Computo
\item Carrera: Ciberseguridad e Infraestructura de Computo
\end{itemize}
\section{Información}
\label{sec:org01ca361}
\subsection{Evaluación:}
\label{sec:org5c10c37}
\begin{itemize}
\item Examenes: 45\%
\item Actividades: 30\%
\item Proyecto final: 25\%
\end{itemize}
\subsection{Herramientas:}
\label{sec:org955c03b}
\begin{itemize}
\item GNS3
\item Packet Tracert
\item Material de Cisco
\end{itemize}
\section{Introducción a la seguridad en redes de computo}
\label{sec:orgafdf53d}
\begin{itemize}
\item Los sistemas de información sostienen la forma en la que la mayoria de las organizaciones hacen negocios.
\item Información: Incluye datos, conocimiento, propiedad intelectual, etc. Que son importantes para una organización.
\item La triada CIA:
\begin{itemize}
\item Confidencialidad
\item Accesibilidad
\item Disponibilidad
\item No repudio: Es la capacidad de demostrar el origen de la creación o modificación de cierta información.
\end{itemize}
\end{itemize}
\subsection{Seguridad de la información}
\label{sec:org7a1fee2}
\begin{itemize}
\item Riesgo: Potencial de perdida, mediante el impacto y la probabilidad de una amenaza aproveche de una vulnerabilidad en uno o mas activos de un sistema de información.
\begin{itemize}
\item Activo: Algo de valor para la organización. Como sistemas de información, información. Pueden ser activos digitales, fisicos e inclusos reputacionales.
\end{itemize}
\item Amenaza: Peligro potencial para los activos, los datos o la red de una organización. Incluye un actor de amenaza o perpetrador de hechos que representan la amenaza.
\begin{itemize}
\item Hackers maliciosos.
\item Empleado deshonestos.
\item Empleado que abre un archivo malicioso.
\item Los actores de amenaza se pueden clasificar en: agentes maliciosos internos, agentes maliciosos externos, internos e internos accidentales.
\end{itemize}
\item Vulnerabilidad: Es la debilidad de un activo (en un sistema o en su diseño) que puede ser explotada por una amenaza.
\begin{itemize}
\item Activos de hardware
\item Activos de software
\end{itemize}
\end{itemize}
\subsubsection{Gesionar el riesgo}
\label{sec:org799c0c2}
\begin{itemize}
\item Proceso de minimizar el potencial de perdida.
\item Implica identificar, evaluar, priorizar riesgos y aplicar medidas de control.
\item Proporcionar transparencia y contabilidad de las operaciones diarias en una organización, alineando los riesgos con los objetivos estrategicos.
\item Ayuda a la organización a prepararse para eventos adversos que podrian impactarle negativamente.
\end{itemize}
\begin{enumerate}
\item Tratamiento del riesgo
\label{sec:org8204893}
\begin{enumerate}
\item Forma de control
\label{sec:orgf2c2f9d}
\begin{itemize}
\item Evitar el riesgo.
\item Reducir el riesgo.
\item Transferir el riesgo.
\item Aceptar el riesgo.
\end{itemize}
\item Controles de mitigación
\label{sec:org97ff614}
\begin{itemize}
\item Metodos por los cuales se intenta proteger cada activo de amenazas (contra-medida que reduce la probabilidad o la gravedad de una amenaza o riesgo).
\begin{itemize}
\item Riesgo residual: Despues de aplicar el control de mitigación, se vuelve un riesgo aceptable.
\end{itemize}
\end{itemize}
\end{enumerate}
\end{enumerate}
\subsubsection{Regulaciones legales y estructuras de cumplimientos}
\label{sec:orgf1e2a31}
\begin{itemize}
\item Los requisitos de la organización podrian proceder de autoridades gubernamentales o de organismos de certificación.
\item Tambien hay estandares como los de la Organización Internacional para la Estandarización (ISO).
\begin{itemize}
\item ISO 27001
\item ISO 27018
\end{itemize}
\item Las organizaciones que cumplen con los requisitos establecidos en la norma estan listas para ser auditadas por terceros y obtener certificaciones o acreditaciones al respecto.
\item Importancia de una certificación:
\begin{itemize}
\item Altamente valoradas por los clientes y empresas asociadas.
\item Evidencia de que se ha alcanzado el nivel minimo de seguridad en la organización.
\item Las regulaciones estan en constante evolución, por lo que la organización debe estar al pendiente de ello.
\end{itemize}
\end{itemize}
\subsubsection{Estrategias de seguridad de la información}
\label{sec:orgf875491}
\begin{itemize}
\item Las organizaciones requieren de programas de seguridad de la información, documentados y actualizados regularmente (la periodicidad de la norma a seguir).
\item Debe incluir:
\begin{itemize}
\item Creación de politicas, procedimientos y manuales de seguridad de la información.
\item Establecimiento y mantenimiento de un programa de concientización.
\item Gestión de riesgos por parte de terceros (servicios en la nube: SLA).
\item Presentación de informes y mejora continua.
\begin{itemize}
\item Captar la atención de personas interesadas.
\item Resaltar la relación costo-beneficio de la estrategia de seguridad de la información.
\end{itemize}
\end{itemize}
\end{itemize}
\subsubsection{Vector de ataque}
\label{sec:orgd8acf1f}
\begin{itemize}
\item Ruta por la que un actor de amenaza puede obtener acceso a un recurso en la red.
\item Se originan dentro y fuera de la red corporativa.
\begin{itemize}
\item Ataque de DoS
\item Robo o copia de datos confidenciales
\item Desconexión de la red.
\end{itemize}
\end{itemize}
\begin{enumerate}
\item Ataques de red habituales
\label{sec:org6c6ef13}
\begin{itemize}
\item Ataques de reconocimiento
\item Ataques de acceso
\item Ingenieria social
\item 
\end{itemize}
\end{enumerate}
\subsection{ACL}
\label{sec:orga8709d3}
\begin{itemize}
\item Proporcionar control de flujo de trafico. Limitar trafico a ciertos enlaces, restringir entrega de actualizaciones.
\end{itemize}
\subsubsection{Funcionamiento}
\label{sec:org4957850}
\begin{itemize}
\item Se pueden aplicar para trafico entrante y saliente
\begin{itemize}
\item Paquetes que ingresan por interfacez de entrada (antes de su busqueda en la table de routing)
\item Paquetes que salen por las interfaces de salida del router (despues de que se enrutan)
\end{itemize}
\item Router compara IPv4 de origen con cada ACE en orden secuencial.
\item Si hay coincidencia se ejecuta la ACE (las ACE restantes )
\end{itemize}
\subsubsection{Mascara Wildcard}
\label{sec:org23c1d05}
\begin{itemize}
\item Para coincidencia con un host
\begin{itemize}
\item Host: 192.168.1.1
\item Wildcard: 0.0.0.0
\end{itemize}

\item Para coincidencia con una subred IPv4
\begin{itemize}
\item 192.168.1.0/24
\item 255.255.255.255 - 255.255.255.0 = 0.0.0.255
\item Mascara Wildcard: 0.0.0.255
\end{itemize}

\item Para coincidencia de distintas subredes
\begin{itemize}
\item 192.168.10.0 y 192.168.11.0
\item Calculas mascara sumarizada: 255.255.254.0
10 = 00001010
11 = 00001011
\item Calculo de mascara wildcard: 0.0.1.255
\end{itemize}

\item Para coincidencia con un rango de direcciones:
\begin{itemize}
\item Ejemplo de rango: 192.168.16.0/24 a 192.168.31.0/24
\item Mascara de subred sumarizada: 255.255.240.0
16 = 00010000
31 = 00011111

\item Mascara wildcard: 255.255.255.255 - 255.255.240.0 = 0.0.15.255
\end{itemize}
\end{itemize}
\begin{enumerate}
\item Palabras clave
\label{sec:org73deaae}
\begin{itemize}
\item Host: Equivale a una mascara wildcard 0.0.0.0
\item Any: Equivale a una mascara wildcard 255.255.255.255
\end{itemize}
\end{enumerate}
\subsection{VPN's}
\label{sec:orgebda1bc}
\begin{itemize}
\item Las organizanciones usan redes privadas virtuales (VPN) para crear conexiones de red privada de extremo a extremo. Una VPN es virtual porque transporta la información dentro de una red privada, pero en realidad, esa información se transporta usando una red publica.
\end{itemize}
\subsection{Tuneles}
\label{sec:orgf69dc39}

\subsection{Tunel PPTP}
\label{sec:orgb6916d4}


\subsection{TLS HANDSHAKE}
\label{sec:orgf5cb008}
\begin{itemize}
\item El intercambio de datos de la aplicacion mediante TLS requiere la negociación previa de un tunel de cifrado.
\item Cliente servidor deben acordar:
\begin{itemize}
\item La versión de TLS a utilizar
\item Elegir el conjunto de cifrado
\item Verificar los certificados
\end{itemize}

\item Se realiza el handshake de TCP
\item Cliente envia versión de TLS, lista de conjunto de cifrado compatibles y otras opciones de TLS que se desee utilizar (CLient Hello)
\item Servidor selecciona version
\end{itemize}
\subsection{RSA}
\label{sec:orgb3062db}
\begin{itemize}
\item Sistema criptográfico de clave publica o asimétrica, de Rivest, Shamir y Adleman.
\item Mecanismo de intercambio de claves predominante en la mayoría de las implementaciones de TLS
\item Esta clave simétrica negociada se utiliza para cifrar su sesión
\end{itemize}
\subsection{Diffie-Hellman}
\label{sec:org1997b29}
\begin{itemize}
\item Cliente y servidor pueden negociar el secreto compartido.

\item Los peers intercambian un valor inicial publicamente.
\item Los peers guardan un valor secreto.
\item Cada peer mezcla el valor publico con el valor secreto, obteniendo un valor distinta cada uno.
\item El valor se intercambia publicamente.
\item Cada peer combina dicho valor con su valor secreto.
\item El resultado es un valor final idéntica en cada peer.
\item El valor final es la clave compartida.
\end{itemize}
\subsection{Open VPN}
\label{sec:orgcfdda45}
\begin{itemize}
\item Es una VPN SSL
\item Implementa la extensión de red segura de capa 2 o 3 mediante el protocolo SSL/TLS
\item Admite metodos de autenticación de cliente
\end{itemize}
\subsection{Configuracion de AutoVPN}
\label{sec:orgfda7d54}
\begin{itemize}
\item apt-get install openvpn easy-rsa
\item ls /etc/openvpn
\item make-cadir easy-rsa (en /etc/openvpn)
\item cd easy-rsa/
\item ls -la
\item ./easyrsa init-pki (en /easy-rsa)
\item ls -la
\item vim /easyrsa/vars
\item ./easyrsa build-ca (en /easy-rsa)
\item poner un password (importante recordar)
\item SRC CA
\item ls pki/ca.crt
\item ls -la pki/
\item ls pki/private
\item cd easy-rsa/
\item ./easyrsa gen-req servidor-src-licic nopass
\item ./easyrsa sign-req server servidor-feiuvi
\item ./easyrsa gen-req cliente1-feiuv nopass
\item ./easyrsa sign-req client cliente1-feiuv
\end{itemize}
\subsection{Detección de intrusiones}
\label{sec:org5b7e839}
\begin{itemize}
\item Supervisar eventos en un sistema informatico o una red y analizarlos en busca de indicios de posibles incidentes.
\item Incidentes: Amenazas inminentes de infracción de las politicas de seguridad informatica. Puede tener multiples causas.

\item IDS: Programa informatico que automatiza el proceso de detección de intrusiones.
\item IPS: Programa informatio que tiene las

\item Enfoque en la identificación de posibles incidentes.
\item No detectan cuando un sistema ya fue comprometido

\item Identificación de problemas en las politicas de seguridad.
\item Documentación de amenazas para la organización.
\item Disuadir a personas de infringir politicas de seguridad.
\end{itemize}
\subsubsection{Funciones clave de los IDPS}
\label{sec:orgf0a37b3}
\begin{itemize}
\item Registro de información sobre los eventos observados
\item Notificación de eventos observados al administrador de seguridad.
\item Generación de reportes.
\end{itemize}
\subsubsection{Tecnicas de respuesta de los IPS}
\label{sec:org0745aec}
\begin{itemize}
\item Detiene el ataque por si mismo.
\item Modifica el entorno de seguridad.
\item Modifica el contenido del ataque.
\end{itemize}
\subsubsection{Precisión en la detección}
\label{sec:org895dd8c}
\begin{itemize}
\item Los IDPS no son totalmente precisos.
\item Falsos positivos: Actividad benigna como maliciosa.
\item Falsos negativos: Actividad maliciosa como benigna.
\end{itemize}
\subsubsection{Metodologias de detección comunes}
\label{sec:orgfa898fb}
\begin{enumerate}
\item Detección basada en firmas
\label{sec:org6d95627}
\begin{itemize}
\item Firma: Un patron que corresponde a una amenaza conocida.
\item Proceso de comparar firmas contra eventos observados para identificar posibles incidentes.
\item Es ineficaz para detectar amenazas no conocidas o camufladas mediante tecnicas de evasión.
\end{itemize}
\item Detección basada en anomalias
\label{sec:orgf07fd1d}
\begin{itemize}
\item Proceso de comparar definiciones de lo que se considera una actividad normal con los eventos observados para identificar desviaciones significativas.
\item Usa perfiles que representan un comportamiento normal de elementos como usuarios, host, conexiones de red o aplicaciones.
\item Es muy efectiva detectando amenazas desconocidas previamente.
\item Periodo de entrenamiento que puede durar dias o semanas.
\end{itemize}
\begin{enumerate}
\item Tipos de perfiles
\label{sec:orgde2da2e}
\begin{itemize}
\item Perfil estatico: No cambia el perfil, solo cuando el IDPS sea configurado para generar un nuevo perfil.
\item Perfil dinamicos: Se ajusta constantemente con eventos adiciones que son observados .
\item Ambientes muy diversos o dinamicos, pueden generar muchos falsos positivos.
\end{itemize}
\end{enumerate}
\item Analisis de protocolos con estado
\label{sec:org4cfb6a2}
\begin{itemize}
\item Proceso de comparar perfiles predeterminados de funciones generalmente aceptadas, como actividades de protocolos benignas para cada estado del protocolo, con los eventos observados y con el fin de identificar desviaciones.
\item Se basa en perfiles universales desarrollados por los provedores que especifican como deben y como no deben utilizarse determinados protocolos.
\item Con estado: Significa que el IDPS es capaz de comprender y rastrear el estado de los protocolos de red, transporte y aplicación (que tienen noción de estado).

\item Protocolo con estado: FTP

\item Permite identificar secuencias de comandos inesperadas
\item Con protocolos que realizan autenticación
\item Comprueba aspectos como la longitud minima y maxima de los argumentos para un comando.
\item Se basan principalmente en estandares de protocolo de proveedores y organismos de normalización como la IETF

\item Consumen muchos recursos, debido a la complejidad del analisis.
\item No pueden detectar ataques que no violen las caracteristicas de comportamiento ``aceptable'' de un protocolo.
\item El modelo de protocolo utilizado por un IDPS podria entrar en conflicto con la forma en que se implementa el protocolo en una aplicación o sistema operativo especifico.
\end{itemize}
\end{enumerate}
\subsubsection{Tipos de tecnologias IDPS}
\label{sec:org7f31daf}
\begin{itemize}
\item Basados en red
\item Inalambricos
\item Basados en analisis de red (NBA)
\item Basados en host2
\end{itemize}
\subsubsection{Componentes tipicos de IDPs}
\label{sec:org768fdaa}
\begin{itemize}
\item Sensor o agente.
\item Servidor de gestión
\item Servidor de gestión
\item Servidor de base de datos
\item Consola
\end{itemize}
\subsubsection{Capacidades de detección}
\label{sec:orgf0b49ae}
\begin{itemize}
\item Umbrales
\item Listas negras: Actividades maliciosas
\item Listas blancas: Actividades benignas
\item Configurar alertas
\item Edición y visualización del codigo
\end{itemize}
\subsubsection{IDPs basado en red}
\label{sec:org999b409}
\begin{itemize}
\item Sensores (appliance)
\item Uno o mas servidores de gestión
\item Multiples consolas
\item Uno o mas servidores de BD (opcional)

\item Las tarjetas estan en modo promiscuo para aceptar todos los paquetes entrantes detectados.

\item Pueden instalarse en uno o dos modos
\begin{itemize}
\item En linea/Inline: El trafico supervisado debe de pasar a traves de el.
\item Pasivo/Passive: Supervisa una copia del trafico real. No pasa a travez del sensor.
\end{itemize}
\end{itemize}
\subsubsection{Capacidades de prevención}
\label{sec:org1605518}
\begin{itemize}
\item Finalizar sesiones TCP
\item Implementar firewall en linea
\item Limite de bandhwith
\item Alteración de contenido malicioso
\item Reconfiguración de dispositivos de red
\item Ejecución de un programa o script
\end{itemize}
\subsubsection{Soluciones de NIDPS Open Source}
\label{sec:org13479d3}
\begin{itemize}
\item Snort
\item Suricata
\item Zeek
\end{itemize}
\end{document}
